% Source-derived chapter 07: Variations of Hodge structure on an analytically general curve
% Original source: geometric-local-systems-general-curves-isomonodromy
\chapter{Variations of Hodge structure on an analytically general curve}
\label{geometric-local-systems-general-curves-isomonodromy-chapter-07}
\source{geometric-local-systems-general-curves-isomonodromy}

\begin{remark}[Source section]
\label{geometric-local-systems-general-curves-isomonodromy-chapter-07-source}
\source{geometric-local-systems-general-curves-isomonodromy}
\source{geometric-local-systems-general-curves-isomonodromy}
This chapter reproduces the corresponding section of the retrieved
LaTeX source, including its definitions, statements, examples, and
proofs. It has not yet been translated into Lean.
\notready
\end{remark}

% The source section title was: Variations of Hodge structure on an analytically general curve
\label{section:hodge-theoretic-results}
\source{geometric-local-systems-general-curves-isomonodromy}
We prove \autoref{theorem:isomonodromic-deformation-CVHS} in
\autoref{subsection:iso-def-proof},
\autoref{theorem:very-general-VHS} in \autoref{subsection:very-general-vhs}, 
\autoref{corollary:abelian-schemes} in
\autoref{subsubsection:geometric-origin-proof},
and
\autoref{corollary:geometric-local-systems} in
\autoref{subsubsection:abelian-scheme-proof}.
Finally, we prove an additional application concerning maps from very general
curves to Hilbert modular varieties in \autoref{subsection:very-general-vhs}.



\subsection{The proof of \autoref{theorem:isomonodromic-deformation-CVHS}}
\label{subsection:iso-def-proof}
\source{geometric-local-systems-general-curves-isomonodromy}
We start with the following lemma, which gives a useful criterion for showing a
monodromy representation is unitary.
\begin{lemma}
\source{geometric-local-systems-general-curves-isomonodromy}
\notready
	\label{lemma:unitary}
\source{geometric-local-systems-general-curves-isomonodromy}
	Suppose $(C, x_1, \ldots, x_n)$ is an $n$-pointed hyperbolic curve and $(E,
	\nabla)$ is a flat vector bundle on $C\setminus\{ x_1, \ldots, x_n\}$
	underlying a polarizable complex variation of Hodge structure. Let
	$(\overline{E}_\star, \overline{\nabla})$ be the Deligne canonical extension of $(E, \nabla)$ to a flat vector bundle on $C$ with regular singularities at the $x_i$, with the parabolic structure associated to $\overline{\nabla}$.
	If $\overline{E}_\star$ is semistable, then the representation of $\pi_1(C\setminus\{x_1, \cdots, x_n\})$ 
	associated to $(E, \nabla)$ is unitary.
\end{lemma}
%We first indicate how \autoref{lemma:unitary} follows from 
%\cite{simpson:iterated-destabilizing-modifications}, and then give a
%separate, self-contained proof.
\begin{remark}
\source{geometric-local-systems-general-curves-isomonodromy}
\notready
	\label{remark:simpson-reference}
\source{geometric-local-systems-general-curves-isomonodromy}
	As pointed out by a referee, \autoref{lemma:unitary} can also be deduced from the results of
	\cite{simpson:iterated-destabilizing-modifications} if $D=\emptyset$, see \cite[\S5.4]{simpson:iterated-destabilizing-modifications}. 
	%Denote the Hodge filtration on $\overline{E}_\star$ by $F$.
	%Using \autoref{prop:basic-facts}, after replacing $(\overline{E}_\star,
	%\overline{\nabla})$ with a direct summand, we can assume
	%$(\overline{E}_\star, \overline{\nabla})$
	%is irreducible, so that the associated Higgs bundle $(\on{gr}_F\overline{E}_\star, \theta)$
	%is stable. 
	%Using 
	%\cite[Theorem 2.5]{simpson:iterated-destabilizing-modifications}, there
	%is a Griffiths-transverse filtration of $\overline{E}_\star$ whose associated graded Higgs bundle semistable. Note that the Hodge filtration $F$ and the trivial filtration both satisfy these conditions.
	 
	%By \cite[Proposition 4.3]{simpson:iterated-destabilizing-modifications},
	%up to shifting the indices on the filtration, this filtration is unique as $(\on{gr}_F\overline{E}_\star, \theta)$ is stable. Hence the Hodge filtration is trivial.
	%Because the Hodge filtration is trivial, the definition of PVHS implies the monodromy of $(\overline{E}_\star, \overline{\nabla})$
	%preserves a definite Hermitian form $Q$ (the polarization), and so the monodromy must factor
	%through the unitary group for $Q$, as desired. 
	Presumably a parabolic version of Simpson's results would imply  \autoref{lemma:unitary} in general, but as far as we know such results have not yet appeared in the literature. 
\end{remark}

\begin{proof}[Proof of \autoref{lemma:unitary}]
	By \autoref{prop:basic-facts}(2), we may write
	$\mathbb{V}:=\ker(\nabla)$ as $$\mathbb{V}:=\bigoplus_i
	\mathbb{L}_i\otimes W_i$$ where the $\mathbb{L}_i$ each have irreducible
	monodromy and also underlie polarizable variations of Hodge structure,
	and the $W_i$ are constant complex Hodge structures.
	It suffices to show the representation associated to each $\mathbb L_i$ has unitary monodromy.
	We may therefore reduce to the case that $\mathbb{V} = \mathbb L_i$ and assume
	that $(E,\nabla)$ has irreducible monodromy.
%Let $(\overline{E}, \overline{\nabla})$ be the Deligne canonical extension of
%$(E, \nabla)$ from $C \setminus \{ x_1, \ldots, x_n\}$ to $C$. 


Let $i$ be maximal such that $F^i\overline{E}_\star$ is non-zero. 
Since $\overline{E}_\star$ is semistable, it follows from
\autoref{corollary:unstable} that
that the natural map 
\begin{align*}
	F^i\overline{E}_\star \to
	F^{i-1}\overline{E}_\star/F^i\overline{E}_\star\otimes \omega_C
\end{align*}
induced by the connection is zero, i.e.~the connection preserves
$F^i\overline{E}_\star$. By irreducibility
of the monodromy of $(E, \nabla)$, we must have that $F^i\overline{E}_\star$ equals
$\overline{E}_\star.$
But in this case $({E}, \nabla)$ is unitary, as the monodromy preserves the polarization, a definite Hermitian form.
\end{proof}


\subsubsection{}
We now recall  the setup of
\autoref{theorem:isomonodromic-deformation-CVHS}. Let
$(C, x_1, \cdots, x_n)$ 
be an $n$-pointed
hyperbolic curve 
of genus $g$.
Let $({E}, \nabla)$ be a flat vector bundle on
$C$ with $\on{rk}{E}<2\sqrt{g+1}$ such that $(E,\nabla)$ has regular singularities at the $x_i$. Our goal is to show that if an isomonodromic
deformation of ${(E, \nabla)}$ to an analytically general nearby $n$-pointed curve underlies a polarizable complex variation of Hodge structure, then $({E},\nabla)$ has unitary monodromy.

\begin{proof}[Proof of \autoref{theorem:isomonodromic-deformation-CVHS}]
As the hypothesis is about the restriction of $(E, \nabla)$ to $C\setminus \{x_1, \cdots, x_n\}$, we may without loss of generality assume $(E,\nabla)$ is the Deligne canonical extension of $(E, \nabla)|_{C\setminus \{x_1, \cdots, x_n\}}$. As local systems underlying complex polarizable variations of Hodge structure are semisimple by \autoref{prop:basic-facts}(2), we may moreover assume that $(E, \nabla)|_{C\setminus \{x_1, \cdots, x_n\}}$ has irreducible monodromy.

Endow $E$ with the parabolic structure associated to $\nabla$, and denote the
corresponding parabolic bundle by $E_\star$.
After replacing $(C, x_1, \cdots, x_n)$ with an analytically general nearby
curve, and $(E,\nabla)$ with its isomonodromic deformation to this curve, we may
assume by \autoref{cor:stable-parabolic} that $E_\star$ is semistable.	Thus $(E, \nabla)$ has unitary monodromy by \autoref{lemma:unitary}.
\end{proof}

\subsection{The proof of \autoref{theorem:very-general-VHS}} 
\label{subsection:very-general-vhs}
\source{geometric-local-systems-general-curves-isomonodromy}
The proof of \autoref{theorem:very-general-VHS} follows  
from the integrality assumption and the following lemma; this lemma is well-known and appears e.g.~in the course of the proofs of \cite[Proposition 4.2.1.3]{katz:pcurvature-and-hodge} or \cite[Proposition 8.2]{esnault2020rigid}, but we have included a proof in the interests of self-containedness.

\begin{lemma}
\source{geometric-local-systems-general-curves-isomonodromy}
\notready
	\label{lemma:integral-and-unitary-implies-finite}
\source{geometric-local-systems-general-curves-isomonodromy}
	Suppose $\Gamma$ is a group,
	%$(C, x_1, \ldots, x_n)$ is an analytically general smooth proper
	%$n$-pointed curve of genus $g$, 
	$K$ is a number field, and $\rho$ is
	a representation $$\rho: \Gamma\to \on{GL}_m(\mathscr{O}_K).$$
	If for each embedding $\iota: K\hookrightarrow \mathbb{C}$ the 
	representation $\rho\otimes_{\mathscr{O}_K, \iota} \mathbb{C}$ is unitary, then
	$\rho$ has finite image.
\end{lemma}
\begin{proof}
Indeed, for $\iota: K\hookrightarrow\mathbb{C}$ an embedding, let $$\rho_\iota:
\Gamma\to \on{GL}_m(\mathbb{C})$$ be the corresponding representation $\rho\otimes_{\mathscr{O}_K, \iota} \mathbb{C}$. First,
$$\prod_\iota \rho_\iota: \Gamma\to \prod_\iota \on{GL}_m(\mathbb{C})$$ has
image contained in a compact set by the definition of being unitary. 
Moreover, the image of $$\mathscr{O}_K\hookrightarrow \prod_\iota \mathbb{C}$$
is discrete, since the difference of any two distinct elements of
$\mathscr{O}_K$ has norm at least $1$.
Hence the image of $\prod_\iota \rho_\iota$ is discrete and has compact closure, and is therefore finite.
\end{proof}


Let $K$ be a number field with ring of integers
$\mathscr{O}_K$. Let $(C, x_1, \cdots, x_n)$ be an analytically general
hyperbolic $n$-pointed curve of genus $g$, and let $\mathbb{V}$ be a $\mathscr{O}_K$-local system on $C\setminus\{x_1, \cdots, x_n\}$ with infinite monodromy.
 Suppose that for each embedding $\iota: \mathscr{O}_K\hookrightarrow
 \mathbb{C}$, $\mathbb{V}\otimes_{\mathscr{O}_K, \iota}\mathbb{C}$ underlies a
 polarizable complex variation of Hodge structure. Our goal is to prove
 \autoref{theorem:very-general-VHS}, which states that $$\on{rk}_{\mathscr{O}_K}(\mathbb{V})\geq 2\sqrt{g+1}.$$

\begin{proof}[Proof of \autoref{theorem:very-general-VHS}]
	We use 
	$\mathscr{T}_{g,n}$ to denote the universal cover of $\mathscr{M}_{g,n}$.
For a fixed representation
$\rho: \pi_1(C\setminus\{x_1, \cdots, x_n\}))\to \on{GL}_m(\mathscr{O}_K)$
let $T_\rho$ denote the set of $[(C', x_1', \cdots, x_n')]\in
\mathscr{T}_{g,n}$,
for which the associated $\mathscr O_K$-local system $\mathbb V$ has the
following property:
for each embedding $\iota: \mathscr{O}_K\hookrightarrow \mathbb{C}$,
$\mathbb{V}\otimes_{\mathscr{O}_K, \iota}\mathbb{C}$ underlies a polarizable
complex variation of Hodge structure on $C' \setminus \{x_1', \ldots, x_n'\}$.
Let $M_\rho$ denote the image of $T_\rho$ under the covering map
$\mathscr{T}_{g,n} \to \mathscr M_{g,n}$.

Our goal is to show that an analytically very general point of $\mathscr M_{g,n}$ lies in the complement of
the union of the $M_\rho,$ where $\rho$ ranges over the set of representations
of $\pi_1(C\setminus\{x_1, \cdots, x_n\})\to \on{GL}_r(\mathscr{O}_K)$, with
infinite image, for $K$ a number field, and $r<2\sqrt{g+1}$.
Since there are only countably many such representations $\rho$,
it is enough to show that an analytically very general point lies in the complement
of $M_\rho$ for fixed $\rho$.
Since $(C, x_1, \ldots, x_n)$ is hyperbolic, 
$\mathscr{T}_{g,n} \to \mathscr M_{g,n}$
is a covering space of countable degree, and so the image of a closed analytic
set is locally contained in a countable union of closed analytic subsets.
It therefore suffices to show that for any $\rho$ with infinite monodromy and
rank $< 2 \sqrt{g+1}$, $T_\rho$ is contained in a closed analytic subset
of $\mathscr T_{g,n}$.

We now show such $T_\rho$ as above are contained in a closed analytic subset of
$\mathscr{T}_{g,n}$.
Indeed, suppose $\mathbb V$ is the local system associated to $\rho$ on some curve $C\setminus\{x_1, \cdots, x_n\}$, and that for each embedding $\iota: \mathscr{O}_K\hookrightarrow \mathbb{C}$, $\mathbb{V}\otimes_{\mathscr{O}_K, \iota} \mathbb{C}$, underlies a polarizable complex variation of Hodge structure. 
It is enough to show this complex polarizable variation of Hodge structure does
not extend to an analytically general nearby curve.
Indeed, if it did,
\autoref{theorem:isomonodromic-deformation-CVHS} implies
$\prod_{\iota: \mathscr O_K \to \mathbb C} \rho_\iota$
has unitary monodromy,
and
\autoref{lemma:integral-and-unitary-implies-finite} implies its monodromy is
finite.
This contradicts our assumption that $\rho$ has infinite monodromy. \end{proof}

\subsection{}\label{subsubsection:geometric-origin-proof}
\source{geometric-local-systems-general-curves-isomonodromy}
\begin{proof}[Proof of \autoref{corollary:geometric-local-systems}]
Let $g\geq 1$ be an integer and let $(C, x_1, \cdots, x_n)$ be an analytically
general hyperbolic $n$-pointed curve of genus $g$. Let $U\subset C\setminus\{x_1, \cdots, x_n\}$ be a dense Zariski-open subset. Let $f: Y\to U$ be a smooth proper morphism, $i\geq 0$ an integer, and suppose $\mathbb{V}$ is a complex local system on $(C, x_1, \cdots, x_n)$ with infinite monodromy such that $\mathbb{V}|_U$ is a summand of $R^if_*\mathbb{C}$. Then we wish to show that $\dim_{\mathbb{C}}\mathbb{V}\geq 2\sqrt{g+1}$.

It suffices to show that $\mathbb{V}$ satisfies the hypotheses of
\autoref{theorem:very-general-VHS}. The existence of an
$\mathscr{O}_K$-structure follows from the fact that $R^if_*\mathbb{C}$ has a
$\mathbb{Z}$-structure. 
Let $\mathbb{W}$ be the corresponding $\mathscr{O}_K$-local system. All that remains is to verify that for each embedding
$\iota:\mathscr{O}_K\hookrightarrow \mathbb{C}$, the corresponding complex local
system $\mathbb{W}_\iota:=\mathbb{W}\otimes_{\mathscr{O}_K, \iota} \mathbb{C}$
underlies a polarizable complex variation of Hodge structure. But  each such
embedding yields a summand $\mathbb{W}_\iota|_U$ of $R^if_*\mathbb{C}$,
Galois-conjugate to the original embedding $\mathbb{W}|_U\subset
\mathbb{V}|_U\subset R^if_*\mathbb{C}.$ Any such summand underlies a polarizable
complex variation of Hodge structure, by \autoref{prop:basic-facts}(2). Now
$\mathbb{W}_\iota|_U$ extends from $U$ to $C\setminus \{x_1, \cdots, x_n\}$ and
so the same is true for the corresponding complex polarizable variation of Hodge
structure by \cite[Corollary 4.11]{schmid1973variation}, which completes the proof.
\end{proof}

\subsection{}
\label{subsubsection:abelian-scheme-proof}
\source{geometric-local-systems-general-curves-isomonodromy}
\begin{proof}[Proof of \autoref{corollary:abelian-schemes}]
	The statement for relative curves $h: S \to C \setminus \{x_1, \ldots, x_n\}$ reduces to the statement for abelian varieties
	upon taking the relative Jacobian $\on{Pic}^0_h$, as the Torelli theorem implies that
	$h$ is isotrivial if $\on{Pic}^0_h$ is.
	Hence, it suffices to show that any abelian scheme $f: A \to C \setminus\{x_1,
	\ldots, x_n\}$ of relative dimension $r < \sqrt{g+1}$ is isotrivial.
	If $r < \sqrt{g+1}$ then $\rk R^1 f_* \mathbb C = 2r < 2 \sqrt{g+1}$, so $R^1 f_* \mathbb C$ has
	finite monodromy by \autoref{corollary:geometric-local-systems}.
	To show $f$ is isotrivial, it is enough to show $f$ is trivial after a finite base change. 
	After passing to a finite \'etale cover of
	$C \setminus\{x_1,\ldots, x_n\}$, we may therefore assume
	$R^1 f_* \mathbb C$ has trivial monodromy.
	In this case, the theorem of the fixed part 
	\cite[Corollaire 4.1.2]{deligne:hodge-ii}
	implies that the first part of the Hodge filtration $F^1(R^1 f_* \mathbb
	C) \subset R^1 f_* \mathbb C$ is also a trivial local system.
	Since $A$ is a quotient of 
	$F^1(R^1 f_* \mathbb C)^\vee$ by the constant local system
	$(R^1 f_* \mathbb Z)^\vee$,
	%$R^{2g-1} f_* \mathbb Z \simeq \wedge^{2g-1} R^1 f_* \mathbb Z$,
	(on fibers this local system is identified with the first homology of $f$
	via Poincar\'e duality,)
	$f$ is also trivial.
\end{proof}


\subsection{}\label{subsubsection:non-density-proof}
\source{geometric-local-systems-general-curves-isomonodromy}
We next prove that local systems of geometric origin with low rank are not
Zariski dense in the character variety of an analytically very general genus $g$
curve.
Recall that we use $\mathscr{M}_{B,r}(X)$ for the character variety
parametrizing conjugacy classes of semisimple representations $$\rho:
\pi_1(C\setminus\{x_1, \cdots, x_n\})\to \on{GL}_r(\mathbb{C}).$$

\begin{proof}[Proof of \autoref{corollary:non-density-of-geometric-local-systems}]
Let $(C, x_1, \cdots, x_n)$ be an analytically very general hyperbolic
$n$-pointed curve of genus $g$, and fix an integer $r$ with $1<r<2\sqrt{g+1}$. 
Note that we must have $g \neq 0$ because there are no integers $r$ with $1 < r
< 2 \sqrt{0 + 1}$.
By \autoref{corollary:geometric-local-systems}, the points of
$\mathscr{M}_{B,r}(C\setminus\{x_1, \cdots, x_n\})$ of geometric origin
correspond to representations $\rho$ with finite image whenever $r < 2
\sqrt{g+1}$. 
We wish to show these finite image representations are not Zariski dense in
$\mathscr{M}_{B,r}(C\setminus\{x_1, \cdots, x_n\})$.
This follows from
by \autoref{lemma:finite-image-not-dense}.
upon choosing a topological identification $C\setminus\{x_1, \cdots, x_n\} \simeq
\Sigma_{g,n}$. 
\end{proof}

\begin{lemma}
\source{geometric-local-systems-general-curves-isomonodromy}
\notready
	\label{lemma:finite-image-not-dense}
\source{geometric-local-systems-general-curves-isomonodromy}
	Let $\Sigma_{g,n}$ be a topological $n$-punctured genus $g$ surface with
	basepoint $p \in \Sigma_{g,n}$. For $r > 1$, the set of representations $\rho: \pi_1(\Sigma_{g,n},p) \to
	\on{GL}_r(\mathbb C)$ with finite image are not Zariski dense in the character
	variety $\mathscr{M}_{B,r}(\Sigma_{g,n}, p)$.
\end{lemma}
\begin{proof}
%	 It suffices to
% prove non-density of representations with finite image in the \emph{framed}
% character variety
% $$\mathscr{M}_{B,r}^\square(\Sigma_{g,n}):=\on{Hom}(\pi_1(\Sigma_{g,n},p),
% \on{GL}_r(\mathbb{C})).$$

 For $s\in \mathbb{Z}_{>0}$, let $V_s\subset
 \mathscr{M}_{B,r}(\pi_1(\Sigma_{g,n},p))$ be the closed subvariety
	 consisting of those representations $\rho$ such that $$\on{Tr} [\rho(x)^s,
	 \rho(y)^s]=\on{Tr}(\on{id})$$ for all $x,y\in \pi_1(\Sigma_{g,n},p)$.

By Jordan's theorem (\cite[p.~91]{jordan1878memoire} or \cite[Theorem
36.13]{curtis1966representation}) on finite subgroups of $\on{GL}_r(\mathbb{C})$,
there is some constant $m_r$ such that for each finite subgroup $G\subset
\on{GL}_r(\mathbb{C})$, there exists an abelian normal subgroup $H\subset G$ of index
at most $m_r$. Hence $V_{(m_r!)}$ contains all representations with finite
image. Thus it suffices to show $V_s$ is not all of
$\mathscr{M}_{B,r}(\Sigma_{g,n})$ for any $s>0$. 

We now write $$\pi_1(\Sigma_{g,n},p)=\left\langle a_1, \cdots, a_g, b_1,\cdots,
b_g, \gamma_1, \cdots, \gamma_n \mid \prod_{i=1}^g [a_i,b_i]\cdot \prod_{j=1}^n
\gamma_j\right\rangle$$ for the standard presentation of the fundamental group.
If $g\geq 2$, let $$\rho: \pi_1(\Sigma_{g,n},p)\to \on{GL}_2(\mathbb{C})$$ be the representation defined by $$a_1 \mapsto \begin{pmatrix} 1 & 1 \\ 0 & 1 \end{pmatrix}$$
$$a_2 \mapsto \begin{pmatrix} 1 & 0 \\ 1 & 1 \end{pmatrix}$$
and sending all other generators to the identity.  If $g=1$, hyperbolicity of
$\Sigma_{g,n}$ implies $n>0$. Then $\pi_1(\Sigma_{g,n},x)$ is free on $n+1$ generators, and we may set $\rho$ to be a representation sending two of the generators to the matrices above, and all other generators to the identity. 
Then $\rho$ does not lie in $V_s$ for any $s>0$ by direct computation. This completes the case $r = 2$. For the case that $r > 2$, the representation $\rho\oplus \on{triv}^{\oplus r-2}$ lies outside of $V_s$ for all $s$, for the same reason.
\end{proof}

\subsection{An application to Hilbert modular varieties}
\label{subsection:hilbert-modular}
\source{geometric-local-systems-general-curves-isomonodromy}

Let $K$ be a totally real number field of degree $h$ over $\mathbb Q$, $\mathscr O_K$ its ring of integers, and 
use $\mathscr X_K$ to denote the Hilbert modular stack parameterizing
$h$-dimensional principally polarized abelian varieties $A$ with an injection
$\mathscr O_K \hookrightarrow \on{End}(A)$.
Using \autoref{corollary:abelian-schemes}, we find there are no nonconstant maps
from an analytically very general $n$-pointed curve of genus $g$
to the moduli stack of principally polarized abelian varieties of dimension $h$
whenever $h \leq \sqrt{g+1}$, or equivalently when $g > h^2$.
We now show that the analogous statement for Hilbert modular stacks holds
whenever $g \geq 1$.
In particular, this improved bound is independent of $h$.
We thank Alexander Petrov for pointing out the following application.

\begin{proposition}
\source{geometric-local-systems-general-curves-isomonodromy}
\notready
	\label{proposition:hilbert-modular}
\source{geometric-local-systems-general-curves-isomonodromy}
	Let $\mathscr X_K$ be the Hilbert modular stack associated to a totally
	real field $K$.
	Let $(C, x_1, \ldots, x_n)$ be an analytically very general hyperbolic $n$-pointed
	curve of genus $g \geq 1$.
	Any map $\phi: C\setminus \{x_1, \ldots, x_n\} \to \mathscr X_K$ is constant.
\end{proposition}
\begin{proof}
	We first use the map $\phi$ to produce a rank $2$ integral local system
	on $C \setminus \{x_1, \ldots, x_n\}$.
	A map $\phi: C\setminus \{x_1, \ldots, x_n\} \to \mathscr X_K$ yields an
	abelian scheme $f: A \to C\setminus \{x_1, \ldots, x_n\}$ with an $\mathscr
	O_K$-action. In particular, $R^1 f_* \mathbb Z$ is a 
	complex PVHS of rank $2h$ with an	$\mathscr O_K$-action.
	If $K$ has degree $h$, then this rank $2h$ local system may not be a
	$\mathscr O_K$-local system, as its fibers may only be locally free, as opposed to free.
	However, after passing to a finite extension $K'$ of $K$, such as the
	Hilbert class field of $K$, any such rank $2$ locally free module
	becomes a free module,
	and so $R^1 f_* \mathbb Z
	\otimes_{\mathscr O_K} \mathscr O_{K'}$ is an $\mathscr O_{K'}$-local system of rank $2$.
	
	We now show the map $\phi$ must be constant.
	By construction,
	for any embedding $\mathscr O_{K'} \to \mathbb C$, the local system
$R^1 f_* \mathbb Z\otimes_{\mathscr O_K} \mathscr O_{K'} \otimes_{\mathscr
O_{K'}} \mathbb C$ underlies a $\mathbb{C}$-PVHS. Hence by \autoref{theorem:very-general-VHS}, $R^1 f_* \mathbb Z\otimes_{\mathscr O_K} \mathscr O_{K'}$ has finite monodromy, as 
\begin{align*}
2 = \rk_{\mathscr O_K'}R^1 f_* \mathbb Z\otimes_{\mathscr O_K} \mathscr O_{K'} < 2 \sqrt{g+1}.
\end{align*}
Hence, the same holds for
$R^1 f_* \mathbb Z$. 
This implies the map $\phi$ is constant using the theorem of the fixed part
\cite[Corollaire 4.1.2]{deligne:hodge-ii} (see the proof
of \autoref{corollary:abelian-schemes} for a similar argument).
\end{proof}
